\documentclass[../../../../SoftwareRequirementsSpecification.tex]{subfiles}
\begin{document}
% Richiede le macro definite in src/macros/useCaseMacros.tex
% (incluse nel preambolo di SoftwareRequirementsSpecification.tex)

\ucstart
\begin{xltabular}{\textwidth}{|l|X|}
  \hline
  \ucid{PT-11} \\ \hline
  \ucname{Cancellazione Scheda} \\ \hline
  \ucactor{PT} \\ \hline
  \ucdescription{Il PT elimina una scheda che ha creato e che non ha
    assegnazioni attive.} \\ \hline
  \ucprecond{Il PT deve essere autenticato. Deve esistere almeno una
    scheda da eliminare. La scheda non deve avere assegnazioni
    attive (vedi Nota).} \\ \hline
  \ucpostcond{Il sistema elimina la scheda con tutti i blocchi
    mensili, le sessioni di allenamento e gli esercizi che
    contiene.} \\ \hline
  \ucmainflow{
    \item Il PT apre la pagina delle proprie schede.
    \item Il sistema mostra l'elenco delle schede create dal PT.
    \item Il PT preme il pulsante "Elimina" su una scheda
      dell'elenco.
    \item Il sistema chiede conferma dell'eliminazione, avvisando
      che l'operazione non è reversibile.
    \item Il PT preme il pulsante "Conferma Eliminazione".
    \item Il sistema verifica che la scheda non abbia assegnazioni
      attive.
    \item Il sistema elimina la scheda.
    \item Il sistema mostra un messaggio di conferma e aggiorna
      l'elenco delle schede.
  } \\ \hline
  \ucaltflow{
    \textbf{4a. Il PT annulla} \newline
    - Il PT preme il pulsante "Annulla" invece di confermare.
    \newline
    - Il sistema non elimina nulla. \newline
    - Il flusso termina. \newline
    \textbf{6a. Scheda con assegnazioni attive} \newline
    - Il sistema mostra un messaggio di errore ("Impossibile
    eliminare: la scheda ha assegnazioni attive") e non elimina la
    scheda. \newline
    - Il flusso termina.
  } \\ \hline
  \ucnote{
    Una scheda con assegnazioni attive non può essere eliminata
    perché uno o più atleti la stanno seguendo: eliminarla
    cancellerebbe le sessioni di allenamento su cui l'atleta si sta
    allenando (vedi caso d'uso A-03). Per smettere di assegnare una
    scheda a un atleta, il PT deve prima terminare l'assegnazione.
    Questa restrizione è la stessa già usata in PT-04 per la
    modifica delle sessioni di allenamento.
  } \\ \hline
\end{xltabular}
\end{document}
